% ============================================================
% Chapter 13 --- Lagrange multipliers
% Originally: content/week-05.tex, \section{Lagrange multipliers} (Corsin Week 5),
%             with problems 5.2, 6.10 and 6.11 from sheets 5 and 6.
% ============================================================
\chapter{Lagrange Multipliers}
\label{ch:lagrange}

% Transition: Claude Opus 5 (Medium)
The previous chapter looked for extrema of $f$ on an open set, where the condition
$\nabla f(x_0) = 0$ is available. Restrict $f$ to a lower-dimensional constraint set and that
condition is generally too strong to hold anywhere; what replaces it is the requirement that
$\nabla f$ be a linear combination of the gradients of the constraints. This chapter develops
that criterion, the geometry behind it, and what goes wrong when the constraint gradients fail to
be independent.

% Originally: content/week-05.tex, \section{Lagrange multipliers} (Corsin Week 5)

% Source: Corsin Nick/Class Notes/Week 5.pdf, pp. 3--4
\section{Lagrange multipliers}
\label{sec:lagrange_multipliers}

Suppose we want to find the hottest place on the surface of the earth. In theory, temperature is
defined everywhere, giving some function
\[T \in C^\infty(\mathbb{R}^3, (0,\infty)) \qquad \text{(absolute zero is not attainable)}\]
If we model the earth as $S^2 := \{x \in \mathbb{R}^3 : |x| = 1\}$, then we want to find the local
maxima of $T|_{S^2}$, which \textbf{are not necessarily critical points of $T$}!

\begin{center}
\begin{tikzpicture}[>=Latex, scale=1.3]
  \draw[blue!80!black, thick] (0,0) circle (1.3);
  \draw[blue!80!black, dashed] (1.3,0) arc (0:180:1.3 and 0.4);
  \draw[blue!80!black] (-1.3,0) arc (180:360:1.3 and 0.4);
  \node[blue!80!black, font=\large] at (-0.6,0.5) {$S^2$};

  % Point q (radial gradient = extremum)
  \coordinate (Q) at (-0.55, 0.95);
  \fill[red!80!black] (Q) circle (1.8pt) node[left, font=\footnotesize] {$q$};
  \draw[red!80!black, thick, ->] (Q) -- ++(-0.4, 0.7) node[above right, font=\small] {$\nabla T(q)$};
  \node[red!80!black, font=\tiny, right] at (-0.1, 1.4) {local extremum of $T|_{S^2}$};

  % Point p (oblique gradient = not extremum)
  \coordinate (P) at (0.35, -0.45);
  \fill[orange!90!black] (P) circle (1.8pt) node[left, font=\footnotesize] {$p$};
  \draw[orange!90!black, thick, ->] (P) -- ++(0.75, -0.2) node[below right, font=\small] {$\nabla T(p)$};
  \draw[orange!90!black, dotted, thick] (-0.3, -0.3) -- (1.0, -0.6);
  \node[orange!90!black, font=\tiny, align=left] at (1.1, -0.1) {has tangent component $\implies$\\not a local extremum};
\end{tikzpicture}
\end{center}

% Generator: Claude Opus 5 (Medium) --- Corsin's sentence reads "T has a local extremum if it is
% perpendicular", which is both missing its subject and the wrong direction of implication.
Read the picture as follows. At $p$ the gradient $\nabla T(p)$ has a non-zero component
\emph{along} the surface, so walking on $S^2$ in that direction increases $T$: the point $p$
cannot be a local extremum of $T|_{S^2}$. So if $T|_{S^2}$ \emph{does} have a local extremum at a
point, then $\nabla T$ there must have no tangential component left --- it must be
\textbf{perpendicular to the surface}, as at $q$.

Note the direction of that statement carefully: perpendicularity is \textbf{necessary, not
sufficient}. Plenty of points have a perpendicular gradient without being extrema at all --- on
the earth, the coldest place has a perpendicular gradient just as the hottest does, and so does
many a saddle. What the condition delivers is a shortlist of \emph{candidates}; deciding which
of them are maxima, which minima, and which neither is a separate step, and in practice is done
by simply evaluating $T$ at each and comparing.

Note also that $S^2$ is the \newterm{level set} (\germanterm{Niveaumenge}) $g(x) = 0$ for
$g(x) := |x|^2 - 1$ (i.e., $S^2 = g^{-1}\{0\}$).

\begin{remark}[Gradient perpendicular to level sets]
For $g(x) := |x|^2-1$ we have $\nabla g(x) = 2x$, which at a point $x \in S^2$ is the outward
radial direction --- perpendicular to $S^2$. In general, the gradient is perpendicular to the
level sets.
\end{remark}

\begin{ainote}
Corsin writes $\nabla g(x) = \tfrac{x}{|x|}$ here. That is the gradient of $x \mapsto |x|$, not
of the $g(x) = |x|^2-1$ just fixed above, whose gradient is $2x$. The two differ only by the
positive factor $2|x|$, so they point the same way and the geometric conclusion is unaffected
--- but the constraint actually differentiated in the Lagrangian below is the squared one. See \cref{ex:xyz_on_ball} for the same $|x|$ versus $|x|^2$ slip in the exercise.
\end{ainote}

% Source: Corsin Nick/Class Notes/Week 5.pdf, pp. 4--5
This intuition is made rigorous in the following.

\begin{proposition}[Constrained optimization problems]
\label{prop:constrained_optimization}
Suppose $f : U \to \mathbb{R}$, $g : U \to \mathbb{R}^k$ are $C^1$ and $U \subseteq \mathbb{R}^n$
is open, and write $g := (g_1,\dots,g_k)$. If $f|_{g^{-1}\{0\}}$ has a local extremum at
$x_0 \in g^{-1}\{0\}$, then
\[\{\nabla f(x_0), \nabla g_1(x_0), \dots, \nabla g_k(x_0)\} \text{ are linearly dependent.}\]
\end{proposition}

\begin{remark}
Once $(\nabla g_1,\dots,\nabla g_k)$ are linearly independent everywhere on $g^{-1}\{0\}$ (the
standing assumption below), the constraint set $g^{-1}\{0\}$ is in fact a submanifold of
$\mathbb{R}^n$ --- this is made precise by \cref{thm:regular_value_theorem}, which also
explains \cref{prop:constrained_optimization} geometrically: $\nabla f(x_0)$ must be normal to the
tangent space of the constraint set at a constrained extremum.
\end{remark}

For $k = 1$, this corresponds to our intuition above of $\nabla f(x_0)$ \textbf{parallel to}
$\nabla g(x_0)$.

\textit{In practice}, we assume $(\nabla g_1, \dots, \nabla g_k)$ nowhere linearly dependent,
since otherwise $\lambda_0 = 0$ is possible (which does not give us anything).

% Generator: Claude Opus 5 (Medium)
\begin{aiexample}[What goes wrong when $\nabla g$ vanishes]
\label{ex:constraint_qualification_fails}
The standing assumption is not bookkeeping --- drop it and the method can miss an extremum
outright. Let
\[f(x,y) := x, \qquad M := \{(x,y) \in \mathbb{R}^2 : y^2 = x^3\},\]
the \newterm{cuspidal cubic} (\germanterm{Neilsche Parabel}). Since $y^2 \geq 0$ forces
$x^3 \geq 0$, every point of $M$ has $x \geq 0$, and $(0,0) \in M$. So $f$ attains a clear
\textbf{global minimum on $M$ at the origin}.

Now run the method with $g(x,y) := y^2-x^3$. The Lagrange condition $\nabla f = \lambda\nabla g$
reads
\[(1,0) = \lambda\,(-3x^2,\ 2y),\]
whose first component demands $1 = -3\lambda x^2$. At the origin this says $1 = 0$: no
multiplier $\lambda$ exists. The method returns \textbf{no candidates at all}, while the minimum
sits there in plain sight.

The culprit is visible in $\nabla g(0,0) = (0,0)$: the standing assumption fails at exactly the
point that matters. Geometrically, $M$ has a \emph{cusp} at the origin --- it is not a smooth
curve there, so it has no tangent line, and \qt{$\nabla f$ perpendicular to the constraint} has
nothing to be perpendicular to.

\textbf{The practical moral:} before trusting a Lagrange computation, check that
$\nabla g \neq 0$ (or, for several constraints, that the $\nabla g_j$ are independent) at every
point of the constraint set. Where it fails, add those points to your candidate list by hand.
\end{aiexample}

% Supplement: Adrien Martelli/Notes/20.03-Notes_Chap11_Lagrangemultipliers_geometric_explanation.pdf
\begin{remark}[Where the coefficient $\lambda_0$ went]
\Cref{prop:constrained_optimization} asserts linear \emph{dependence} of
$\{\nabla f(x_0), \nabla g_1(x_0), \dots, \nabla g_k(x_0)\}$, i.e.\ scalars
$\lambda_0, \lambda_1, \dots, \lambda_k$, not all zero, with
$\lambda_0 \nabla f(x_0) + \lambda_1 \nabla g_1(x_0) + \dots + \lambda_k \nabla g_k(x_0) = 0$.
The standing assumption above --- $(\nabla g_1,\dots,\nabla g_k)$ linearly independent --- is
exactly what rules out $\lambda_0 = 0$ (else the $\nabla g_j$'s alone would already be dependent),
so $\lambda_0 \neq 0$ and we may divide by it, silently renaming $\lambda_j/\lambda_0$ as
$\lambda_j$; this is why the Lagrangian below carries no separate $\lambda_0$. Adrien Martelli's
notes keep $\lambda_0$ explicit instead of dividing it away, pairing the dependence relation with
a normalization constraint $\lambda_0^2 + \dots + \lambda_k^2 = 1$ (any nonzero vector of
multipliers can be rescaled to satisfy this). The two formulations are equivalent whenever
$\lambda_0 \neq 0$; keeping $\lambda_0$ explicit only matters in degenerate settings where the
independence assumption on $(\nabla g_1,\dots,\nabla g_k)$ fails and $\lambda_0 = 0$ is genuinely
forced, e.g.\ if $x_0$ is itself a critical point of the constraint map $g$.
\end{remark}

% Generator: Claude Opus 5 (High)
\begin{remark}[What the multiplier itself measures]
\label{rem:lagrange_multiplier_sensitivity}
So far $\lambda$ has been bookkeeping --- a number produced by the method and then discarded once
the candidate points are found. It carries information of its own. Consider the constrained
problem with the constraint level left free,
\[m(c) := \min\{f(x) : g(x) = c\},\]
and suppose the minimiser $x^*(c)$ depends differentiably on $c$. Differentiating
$g(x^*(c)) = c$ gives $\langle \nabla g(x^*(c)), (x^*)'(c)\rangle = 1$, while the Lagrange
condition $\nabla f = \lambda\nabla g$ at $x^*(c)$ turns the chain rule for $m$ into
\[m'(c) = \big\langle \nabla f(x^*(c)),\ (x^*)'(c)\big\rangle
        = \lambda\big\langle \nabla g(x^*(c)),\ (x^*)'(c)\big\rangle = \lambda.\]
So $\lambda$ is the rate at which the optimal value responds to relaxing the constraint by one
unit --- in economics this reading is standard enough to have its own name, the
\qt{shadow price} of the constraint. A large $\lvert\lambda\rvert$ says the constraint is
expensive; $\lambda = 0$ says it is not binding at all.
\end{remark}

% Source: Corsin Nick/Class Notes/Week 5.pdf, pp. 4--5
% Supplement: Sascha Brack/Class Notes/Week_05_Notes_Friday.pdf, p. 5
\begin{center}
\begin{tikzpicture}[>=Latex, scale=1.3]
  \draw[very thick] (0,0) circle (1.2);
  \node[font=\large] at (-0.5,0.7) {$M$};
  \fill (0,0) circle (1.2pt) node[below right, font=\scriptsize] {$0$};

  \coordinate (P1) at ({1.2*cos(20)}, {1.2*sin(20)});
  \draw[magenta!90!black, thick, ->] (P1) -- ++({0.8*cos(20)}, {0.8*sin(20)}) node[above, font=\scriptsize] {$\nabla g(x_1, y_1)$};
  \fill (P1) circle (1.5pt) node[below left, font=\scriptsize, magenta!90!black] {$(x_1, y_1)$};

  \coordinate (P2) at ({1.2*cos(60)}, {1.2*sin(60)});
  \draw[magenta!90!black, thick, ->] (P2) -- ++({0.8*cos(60)}, {0.8*sin(60)}) node[right, font=\scriptsize] {$\nabla g(x_2, y_2)$};
  \fill (P2) circle (1.5pt) node[below left, font=\scriptsize, magenta!90!black] {$(x_2, y_2)$};

  \coordinate (P3) at ({1.2*cos(150)}, {1.2*sin(150)});
  \draw[magenta!90!black, thick, ->] (P3) -- ++({0.8*cos(150)}, {0.8*sin(150)}) node[above left, font=\scriptsize] {$\nabla g(x_3, y_3)$};
  \fill (P3) circle (1.5pt) node[below right, font=\scriptsize, magenta!90!black] {$(x_3, y_3)$};
\end{tikzpicture}
\end{center}

As the figure suggests, $\nabla g(x)$ is perpendicular to the constraint set $M := g^{-1}\{0\}$. To find extrema along $M$, we basically only care about the tangential part of $\nabla f$. We can then find a multiplier $\lambda$ to \qt{erase} the perpendicular component of $\nabla f(x_0)$. We do this by forming the \newterm{Lagrangian}
(\germanterm{Lagrange-Funktion}):
\[
\begin{aligned}
L : U \times \mathbb{R}^k &\to \mathbb{R} \\
(x,\lambda) &\mapsto f(x) - \lambda\cdot g(x) = f(x) - \lambda_1g_1(x) - \dots - \lambda_kg_k(x)
\end{aligned}
\]
We then find local extrema by forcing $\nabla L(x,\lambda) = 0$, i.e.
\[\underbrace{\frac{\partial L}{\partial x_i} = 0 \ \ \forall i = 1,\dots,n}_{\nabla f(x) - \lambda\cdot\nabla g(x) = 0} \qquad \text{and} \qquad \underbrace{\frac{\partial L}{\partial \lambda_j} = 0 \ \ \forall j = 1,\dots,k}_{g_j(x) = 0}\]

% Generator: Claude Opus 5 (Medium)
\begin{remark}[Why the Lagrangian is built this way]
\label{rem:why_the_lagrangian}
The construction can look like a trick pulled out of nowhere. It is not: $L$ is assembled
precisely so that setting \emph{one} gradient to zero encodes \emph{both} halves of the problem
at once.

Read the two underbraces. Differentiating in the $x_i$ reproduces the geometric condition ---
$\nabla f$ parallel to the $\nabla g_j$, i.e.\ no tangential component left. Differentiating in
$\lambda_j$ returns $g_j(x)=0$, the constraint itself, because $\lambda_j$ appears in $L$ only
in the term $-\lambda_jg_j(x)$, linearly. So the multipliers are not auxiliary clutter: they are
exactly the bookkeeping that lets a \emph{constrained} problem in $n$ variables be written as an
\emph{unconstrained} critical-point problem in $n+k$ variables, which is the one kind of problem
we already know how to solve (\cref{prop:extremum_implies_critical}).

That also explains the minus sign, which is arbitrary: writing $f + \lambda\cdot g$ merely flips
the sign of every $\lambda_j$ and changes nothing. And it explains why $\lambda$ is discarded at
the end --- it is scaffolding, not part of the answer.

\begin{ainote}
Though $\lambda$ is not meaningless. In applications it is the \newterm{shadow price}
(\germanterm{Schattenpreis}): if the constraint is relaxed from $g=0$ to $g=c$, the optimal value
changes at rate $\lambda$ to first order, $\tfrac{d}{dc}f_{\text{opt}}(c) = \lambda$. This is
where the name comes from in economics, and it is why $\lambda$ carries units. Not needed for
this course, but it is the usual answer to \qt{what \emph{is} $\lambda$?}
\end{ainote}
\end{remark}

% Source: Corsin Nick/Class Notes/Week 5.pdf, pp. 5--6
\begin{exercise}[Extrema of $xyz$ on the unit ball]
\label{ex:xyz_on_ball}
Find the local extrema of
\[f : K \to \mathbb{R}, \quad (x,y,z) \mapsto xyz, \qquad K = \{(x,y,z) \in \mathbb{R}^3 : x^2+y^2+z^2 \leq 1\}.\]
\end{exercise}

% Kept inline: solved directly in Corsin Nick/Class Notes/Week 5.pdf, pp. 5--6.
\begin{exercisesolution}[Proof of \cref{ex:xyz_on_ball}]
We first find the \textbf{interior critical points}:
\[\nabla f(x,y,z) = (yz,\ xz,\ xy) \overset{!}{=} 0 \implies x = y = 0 \ \text{ or } \ x = z = 0 \ \text{ or } \ y = z = 0.\]
Then we look for local extrema \textbf{on the boundary}. We need to find the local extrema of
$f|_{g^{-1}\{0\}}$, where $g : \mathbb{R}^3\to\mathbb{R}$ is given by $g(x) = |x|^2 - 1$.

The Lagrangian is
\[L(x,y,z,\lambda) := f(x,y,z) - \lambda g(x,y,z) = xyz - \lambda(x^2+y^2+z^2-1),\]
giving the system of equations
\[
\begin{aligned}
\frac{\partial L}{\partial x} &= yz - 2\lambda x = 0, & \frac{\partial L}{\partial \lambda} &= 1 - (x^2+y^2+z^2) = 0. \\
\frac{\partial L}{\partial y} &= xz - 2\lambda y = 0, \\
\frac{\partial L}{\partial z} &= xy - 2\lambda z = 0,
\end{aligned}
\]
We then have
\begin{equation*}
2\lambda = \frac{yz}{x} = \frac{xz}{y} = \frac{xy}{z} \tag{$*$}
\end{equation*}
\[\implies \underbrace{y^2z^2 = x^2z^2}_{y^2 = x^2} = \underbrace{x^2y^2}_{y^2 = z^2} \qquad \text{(if any coordinate is 0, so are the others)}\]
\[\implies x^2 = y^2 = z^2.\]
Then $x^2+y^2+z^2 = 1 \implies 3x^2 = 1 \implies x = \pm\frac{1}{\sqrt{3}}$, giving
\textbf{8 solutions}:
\[(x,y,z) \in \left\{\tfrac{1}{\sqrt{3}}(\alpha,\beta,\gamma) : \alpha,\beta,\gamma \in \{\pm 1\}\right\}\]
plus the solutions $\{(x_1,x_2,x_3) \in K : x_i = x_j = 0,\ i \neq j\}$.

\textbf{To justify division by $z$, $x$ or $y$ in $(*)$:} if $z = 0$, then
$-\lambda x = -\lambda y = xy = 0$, giving $x = 0$ or $y = 0$, and $\lambda = 0$, since
$x = y = z = 0$ does not satisfy $x^2+y^2+z^2 = 1$. Then, since $\lambda = 0$,
$\nabla f(x,y,z) = 0$ and we have already found that point in the first step. Similarly for
$y = 0$ or $x = 0$.

% Generator: Claude Opus 5 (Medium) --- Corsin's solution stops at the candidate list. The
% exercise asks for the local extrema, so the classification step is completed here.
\textbf{Classifying the candidates.} The work so far produced candidates, not answers
(\cref{prop:extremum_implies_critical} and the warning after it). Since $K$ is compact and $f$
is continuous, $f$ does attain a global maximum and minimum on $K$, so it is enough to evaluate.

At a boundary candidate $\tfrac{1}{\sqrt3}(\alpha,\beta,\gamma)$ with
$\alpha,\beta,\gamma \in \{\pm1\}$,
\[f\Big(\tfrac{1}{\sqrt3}(\alpha,\beta,\gamma)\Big)
  = \frac{\alpha\beta\gamma}{3\sqrt3},\]
so the value depends only on the sign $\alpha\beta\gamma$. Four of the eight points have
$\alpha\beta\gamma = +1$ and give the maximal value $\tfrac{1}{3\sqrt3}$; the other four have
$\alpha\beta\gamma = -1$ and give $-\tfrac{1}{3\sqrt3}$. On the coordinate axes found in the
first step, $f \equiv 0$.

Hence $f$ attains its \textbf{global maximum} $\tfrac{1}{3\sqrt3}$ at the four points with an
even number of minus signs, and its \textbf{global minimum} $-\tfrac{1}{3\sqrt3}$ at the four
with an odd number. The interior candidates are \textbf{not extrema at all}: at any point of a
coordinate axis $f=0$, yet every neighbourhood contains points where $xyz>0$ and points where
$xyz<0$ --- flip the sign of one coordinate. They are the three-variable analogue of a saddle.
\end{exercisesolution}

\begin{ainote}
Corsin writes $g(x) = |x| - 1$ for the constraint above, but the Lagrangian uses
$x^2+y^2+z^2-1$, i.e.\ $|x|^2-1$. Both cut out the same sphere; the squared version is the one
actually differentiated, so it is used throughout.
\end{ainote}

% Source: Jérôme Paschoud/Notizen/Woche 5.2 Optimierung I.pdf, p. 3
% Generator: Gemini 3.6 Flash (High)
\begin{exercise}[Extrema on a circle]
\label{ex:extrema_on_circle}
Find the maximum and minimum of $f(x,y) = xy$ under the condition that $(x,y)$ lies on the circle with radius $\sqrt{2}$. (cf.\ Michaels Bsp 7.2.3)
\end{exercise}

\begin{exercisesolution}[\cref{ex:extrema_on_circle}]
We want to optimize $f(x,y) = xy$ subject to the constraint $g(x,y) = x^2+y^2-2 = 0$.
The Lagrangian is
\[L(x,y,\lambda) := f(x,y) - \lambda g(x,y) = xy - \lambda(x^2+y^2-2).\]
Setting the gradient of $L$ to zero yields the system:
\begin{align*}
  \partial_x L &= y - 2\lambda x = 0 \\
  \partial_y L &= x - 2\lambda y = 0 \\
  \partial_\lambda L &= 2 - (x^2+y^2) = 0.
\end{align*}
From the first two equations, $y = 2\lambda x$ and $x = 2\lambda y$. Substituting the second into the first gives $y = 4\lambda^2 y$, or $y(1-4\lambda^2) = 0$.
If $y=0$, then $x=0$, which contradicts $x^2+y^2=2$. Hence $y \neq 0$ and $4\lambda^2 = 1 \implies \lambda = \pm 1/2$.
If $\lambda = 1/2$, then $y = x$. The constraint gives $2x^2 = 2 \implies x = \pm 1$, giving points $(1,1)$ and $(-1,-1)$. At these points, $f(\pm 1, \pm 1) = 1$.
If $\lambda = -1/2$, then $y = -x$. The constraint gives $2x^2 = 2 \implies x = \pm 1$, giving points $(1,-1)$ and $(-1,1)$. At these points, $f(\pm 1, \mp 1) = -1$.
Since the circle is compact, these must be the global extrema. The maximum is $1$ and the minimum is $-1$.
\end{exercisesolution}

% Source: Jérôme Paschoud/Notizen/Woche 5.2 Optimierung I.pdf, p. 3
% Generator: Gemini 3.6 Flash (High)
\begin{exercise}[Largest rectangle in an ellipse]
\label{ex:largest_rectangle_ellipse}
Given an ellipse $E = \{ (x,y) \in \mathbb{R}^2 \mid x^2/A^2 + y^2/B^2 = 1 \}$, find the rectangle (with sides parallel to the axes) whose vertices lie on $E$ and has the largest area. (cf.\ Michaels Bsp 7.2.7)
\end{exercise}

\begin{exercisesolution}[\cref{ex:largest_rectangle_ellipse}]
Let the vertices of the rectangle in the four quadrants be $(\pm x, \pm y)$ with $x, y \geq 0$. The area of the rectangle is $f(x,y) = 4xy$.
We maximize $f(x,y)$ subject to the constraint $g(x,y) = \frac{x^2}{A^2} + \frac{y^2}{B^2} - 1 = 0$ in the region $x > 0, y > 0$.
The Lagrangian is
\[L(x,y,\lambda) := 4xy - \lambda\left(\frac{x^2}{A^2} + \frac{y^2}{B^2} - 1\right).\]
The critical point equations are:
\begin{align*}
  \partial_x L &= 4y - \frac{2\lambda x}{A^2} = 0 \implies \lambda = \frac{2yA^2}{x} \\
  \partial_y L &= 4x - \frac{2\lambda y}{B^2} = 0 \implies \lambda = \frac{2xB^2}{y} \\
  \partial_\lambda L &= 1 - \frac{x^2}{A^2} - \frac{y^2}{B^2} = 0.
\end{align*}
Equating the expressions for $\lambda$ yields $\frac{2yA^2}{x} = \frac{2xB^2}{y} \implies y^2 A^2 = x^2 B^2$, or $\frac{x^2}{A^2} = \frac{y^2}{B^2}$.
Substituting this into the constraint gives $2\frac{x^2}{A^2} = 1 \implies x = \frac{A}{\sqrt{2}}$.
Similarly, $y = \frac{B}{\sqrt{2}}$.
The maximal area is $f\left(\frac{A}{\sqrt{2}}, \frac{B}{\sqrt{2}}\right) = 4 \cdot \frac{A}{\sqrt{2}} \cdot \frac{B}{\sqrt{2}} = 2AB$.
\end{exercisesolution}

% Source: Corsin Nick/Class Notes/Week 5.pdf, p. 6
\begin{center}
\begin{tikzpicture}[>=Latex, scale=1.4]
  % Sphere
  \draw[blue!80!black, thick] (0,0) circle (1.4);
  \draw[blue!80!black, dashed] (1.4,0) arc (0:180:1.4 and 0.45);
  \draw[blue!80!black] (-1.4,0) arc (180:360:1.4 and 0.45);

  % Axes
  \draw[green!60!black, thick, ->] (-1.6,0) -- (1.6,0) node[right, font=\footnotesize] {$x$};
  \draw[green!60!black, thick, ->] (0,-1.6) -- (0,1.6) node[above, font=\footnotesize] {$z$};
  \draw[green!60!black, thick, ->] (0.8,-1.0) -- (-0.8,1.0) node[above left, font=\footnotesize] {$y$};

  % Inscribed cube vertices (+-1/sqrt(3), +-1/sqrt(3), +-1/sqrt(3))
  \coordinate (A1) at (-0.6, -0.5);
  \coordinate (A2) at (0.6, -0.5);
  \coordinate (A3) at (0.8, 0.2);
  \coordinate (A4) at (-0.4, 0.2);

  \coordinate (B1) at (-0.6, 0.7);
  \coordinate (B2) at (0.6, 0.7);
  \coordinate (B3) at (0.8, 1.4);
  \coordinate (B4) at (-0.4, 1.4);

  % Cube edges
  \draw[red!80!black, thick, dotted] (A1)--(A2)--(A3)--(A4)--cycle;
  \draw[red!80!black, thick, dotted] (B1)--(B2)--(B3)--(B4)--cycle;
  \draw[red!80!black, thick, dotted] (A1)--(B1) (A2)--(B2) (A3)--(B3) (A4)--(B4);

  % Vertex dots
  \foreach \p in {A1,A2,A3,A4,B1,B2,B3,B4} {
    \fill[red!80!black] (\p) circle (1.8pt);
  }
  \node[below, font=\small] at (0,-1.7) {8 critical points $\tfrac{1}{\sqrt{3}}(\pm 1, \pm 1, \pm 1)$ on $S^2$};
\end{tikzpicture}
\end{center}

% Generator: Claude Opus 5 (High) --- not in Corsin's notes; included because it is the one
% place in the course where an analytic tool repays a debt to linear algebra, and because the
% Rayleigh-quotient argument is a genuine use of \cref{prop:constrained_optimization} with more
% than one constraint.
\begin{theorem}[Spectral theorem, via constrained optimization]
\label{thm:spectral_theorem_optimization}
Every symmetric $A \in \mathbb{R}^{n\times n}$ admits an orthonormal basis
$\{v_1,\dots,v_n\}$ of $\mathbb{R}^n$ consisting of eigenvectors of $A$, with real eigenvalues.
\end{theorem}

\begin{proof}
Let $f(x) := \langle x, Ax\rangle$, a polynomial and hence continuous, and let
$g(x) := \lvert x\rvert^2 - 1$, so that the unit sphere is $S^{n-1} = g^{-1}\{0\}$. Since $A$ is
symmetric, $\nabla f(x) = 2Ax$; also $\nabla g(x) = 2x$.

We construct the $v_j$ one at a time. Suppose $v_1,\dots,v_k$ are already orthonormal
eigenvectors of $A$, with eigenvalues $\lambda_1,\dots,\lambda_k$ (for $k=0$ there is nothing to
suppose). Put
\[K_k := S^{n-1} \cap \{x : \langle x,v_j\rangle = 0 \text{ for } j=1,\dots,k\},\]
the unit sphere of the orthogonal complement of $\operatorname{span}\{v_1,\dots,v_k\}$. This is
closed and bounded, hence compact by \cref{thm:heine_borel}, and non-empty as long as $k<n$, so
$f$ attains a maximum on it at some $v_{k+1} \in K_k$.

The constraints active at $v_{k+1}$ are $g$ together with the $k$ linear functions
$h_j(x) := \langle x,v_j\rangle$, whose gradients are $\nabla g(v_{k+1}) = 2v_{k+1}$ and
$\nabla h_j(v_{k+1}) = v_j$. These $k+1$ vectors are orthogonal and non-zero, hence linearly
independent, so \cref{prop:constrained_optimization} applies and yields multipliers with
\[2Av_{k+1} = 2\lambda_{k+1}v_{k+1} + \sum_{j=1}^{k}\mu_jv_j.\]
Pair both sides with $v_i$ for a fixed $i \leq k$. On the right, $\langle v_{k+1},v_i\rangle = 0$
and $\langle v_j,v_i\rangle = \delta_{ji}$, leaving $\mu_i$. On the left, symmetry of $A$ moves it
across the inner product,
\[\langle Av_{k+1}, v_i\rangle = \langle v_{k+1}, Av_i\rangle
  = \lambda_i\langle v_{k+1},v_i\rangle = 0,\]
using that $v_i$ is already an eigenvector. Hence $\mu_i = 0$ for every $i \leq k$, and the
displayed identity collapses to $Av_{k+1} = \lambda_{k+1}v_{k+1}$. So $v_{k+1}$ is a unit
eigenvector orthogonal to all its predecessors, and after $n$ steps the $v_j$ form the required
basis. Each eigenvalue is real because $\lambda_j = \langle v_j, Av_j\rangle$ is a real number by
construction.
\end{proof}

% Generator: Claude Opus 5 (High)
\begin{corollary}[Extreme eigenvalues as Rayleigh quotients]
\label{cor:courant_fischer_extreme}
For symmetric $A \in \mathbb{R}^{n\times n}$,
\[\lambda_{\min}(A) = \min_{x \in S^{n-1}}\langle x,Ax\rangle, \qquad
  \lambda_{\max}(A) = \max_{x \in S^{n-1}}\langle x,Ax\rangle.\]
\end{corollary}

\begin{proof}
Write $x \in S^{n-1}$ in the orthonormal eigenbasis of \cref{thm:spectral_theorem_optimization}
as $x = \sum_j c_jv_j$ with $\sum_j c_j^2 = 1$. Then
$\langle x,Ax\rangle = \sum_j \lambda_jc_j^2$ is a convex combination of the eigenvalues, so it
lies between $\lambda_{\min}$ and $\lambda_{\max}$; both bounds are attained by taking $x$ to be
the corresponding eigenvector.
\end{proof}

% Originally: content/week-05.tex, end of the chapter

\begin{aiexercise}[Extrema on a Circle via Lagrange Multipliers]
\label{ex:ai_lagrange_circle}
% Generator: Gemini 3.6 Flash (Medium)
Find the maximum and minimum values of $f(x,y) := x + 2y$ subject to the constraint $g(x,y) := x^2 + y^2 - 5 = 0$.
\end{aiexercise}


% --- exercises relocated here from the dissolved sheets ---

% Originally: content/week-05.tex, \exercisesheet{5}, problem 5.2
% Source: exercises/Ex5_Analysis2_eng.pdf

\begin{exercise}[5.2 --- Lagrange multipliers \textnormal{(important)}]
\label{ex:5.2}
Consider the ball $U := \{(x,y,z) : x^2+y^2+z^2 < 1\}$ and the function
\[f(x,y,z) := 1 - x^2 - y^2 + 4x.\]
\begin{enumerate}[label=\textbf{\arabic*.}]
  \item Motivate rigorously why $f$ attains its absolute maximum and minimum in $\overline{U}$.
  \item Find all the critical points of $f$ which lie in the interior of $U$. Say whether
    $\sup_U f$ (or $\inf_U f$) is attained at some point in $U$.
  \item Say whether $\max_{\partial U} f$ and $\min_{\partial U}$ exist.
  \item With the method of Lagrange multipliers, find the possible critical points for the
    constrained problem $\max\{f(x,y,z) : (x,y,z) \in \partial U\}$, and same for min.
  \item Among all the points you have found, clarify at which points the maximum/minimum of $f$
    is attained.
\end{enumerate}
\end{exercise}

% Originally: content/week-06.tex, \exercisesheet{6}, problem 6.10
% Source: exercises/Ex6_Analysis2_eng.pdf

\begin{exercise}[6.10 --- Lagrange multipliers \textnormal{(important)}]
\label{ex:6.10}
Consider the function $f(x,y,z) = 3x-y+2z$ and the set
\[M = \{(x,y,z) \in \mathbb{R}^3 \mid x^2+y^2+z^2 = 1,\ x+y = 0\}.\]
Determine the extrema of $f$ on $M$ and their nature.
\end{exercise}

% Originally: content/week-06.tex, \exercisesheet{6}, problem 6.11
% Source: exercises/Ex6_Analysis2_eng.pdf

\begin{exercise}[6.11 --- Closest point on a hyperbolic paraboloid \textnormal{(important)}]
\label{ex:6.11}
Find the points on the submanifold
\[M = \{(x,y,z) \in \mathbb{R}^3 : z = x^2-y^2\}\]
that are closest to the point $(0,0,1)$.
\end{exercise}

% Originally: the \section{Solutions} of content/week-05.tex and content/week-06.tex

\section{Solutions}
\label{sec:lagrange_solutions}

\begin{exercisesolution}[Solution to \cref{ex:5.2}]
% Generator: Claude Sonnet 5 (Medium)
\begin{enumerate}[label=\textbf{\arabic*.}]
  \item $\overline{U}$ is closed and bounded in $\mathbb{R}^3$, hence compact by
    \Cref{thm:heine_borel}, and $f$ is a polynomial, hence continuous. By the extreme value
    theorem (\cref{item:extreme_value_theorem}), $f$ attains its absolute maximum and minimum on
    $\overline{U}$.
  \item $\nabla f(x,y,z) = (-2x+4,\ -2y,\ 0)$ vanishes only at $x=2$, $y=0$ (any $z$), and no such
    point lies in $U$ (since $x=2$ already violates $x^2+y^2+z^2<1$). So $f$ has \textbf{no
    critical point in the interior of $U$}; consequently, by
    \Cref{prop:extremum_implies_critical}, $\sup_U f$ and $\inf_U f$ are \textbf{not attained} at
    any interior point.
  \item $\partial U$ is compact (a sphere) and $f$ is continuous, so $\max_{\partial U} f$ and
    $\min_{\partial U} f$ \textbf{both exist}, again by the extreme value theorem
    (\cref{item:extreme_value_theorem}).
  \item With $g(x,y,z) := x^2+y^2+z^2-1$, the Lagrange condition $\nabla f = \lambda \nabla g$
    reads $(-2x+4,-2y,0) = \lambda(2x,2y,2z)$. The last coordinate gives $\lambda z = 0$.
    \textbf{If $z \neq 0$:} then $\lambda = 0$, so $-2x+4=0 \Rightarrow x=2$, impossible on
    $\partial U$. So $z = 0$. Then $x^2+y^2=1$, and the remaining equations are
    $4-2x = 2\lambda x$ and $-2y = 2\lambda y$, i.e.\ $-2y(1+\lambda) = 0$.
    \textbf{If $\lambda = -1$:} the first equation becomes $4-2x=-2x$, i.e.\ $4=0$, impossible.
    So $y = 0$, and $x^2=1$ gives $x = \pm 1$. The Lagrange candidates are
    $(1,0,0)$ (with $\lambda = 1$) and $(-1,0,0)$ (with $\lambda = -3$).
  \item Evaluating, $f(1,0,0) = 1-1+4 = 4$ and $f(-1,0,0) = 1-1-4 = -4$. Since there are no
    interior critical points, these are the global extrema on $\overline{U}$: the
    \textbf{maximum} $4$ is attained at $(1,0,0)$, and the \textbf{minimum} $-4$ at $(-1,0,0)$.
\end{enumerate}
\end{exercisesolution}

\begin{exercisesolution}[Proof of \cref{ex:ai_lagrange_circle}]
% Generator: Gemini 3.6 Flash (Medium)
The constraint set $S := \{(x,y) \in \mathbb{R}^2 \mid x^2 + y^2 = 5\}$ is compact and $\nabla g(x,y) = (2x, 2y)\transp \neq (0,0)\transp$ on $S$. By Lagrange Multipliers, candidates satisfy $\nabla f = \lambda \nabla g$, i.e.:
\[
1 = 2\lambda x \quad \text{and} \quad 2 = 2\lambda y \implies y = 2x.
\]
Substituting $y = 2x$ into $x^2 + y^2 = 5$ gives $5x^2 = 5 \implies x = \pm 1$. The candidate points are $(1,2)$ and $(-1,-2)$. Evaluating $f$:
\[
f(1,2) = 1 + 4 = 5 \quad (\text{maximum}), \qquad f(-1,-2) = -1 - 4 = -5 \quad (\text{minimum}).
\]
\end{exercisesolution}

\begin{exercisesolution}[Solution to \cref{ex:6.10}]
% Generator: Claude Sonnet 5 (Medium)
With $g_1(x,y,z) := x^2+y^2+z^2-1$ and $g_2(x,y,z) := x+y$, the Lagrangian
$L := f - \lambda_1 g_1 - \lambda_2 g_2$ gives the system
\[3 - 2\lambda_1 x - \lambda_2 = 0, \qquad -1-2\lambda_1 y - \lambda_2 = 0, \qquad 2 - 2\lambda_1 z = 0, \qquad x+y=0, \qquad x^2+y^2+z^2=1.\]
From $x+y=0$, $y=-x$. Subtracting the second equation from the first eliminates $\lambda_2$:
\[4 - 2\lambda_1 x - 2\lambda_1(-x) \cdot(-1)= 4-2\lambda_1x-2\lambda_1x=4-4\lambda_1 x = 0 \implies \lambda_1 x = 1.\]
From $2-2\lambda_1 z=0$, $\lambda_1 = 1/z$ (so $z \neq 0$), hence $x = z$ (from $\lambda_1 x = 1 =
\lambda_1 z$). With $y=-x=-z$, the constraint becomes $3x^2 = 1$, i.e.\ $x = \pm\tfrac{1}{\sqrt3}$,
giving the two points $\tfrac{1}{\sqrt3}(1,-1,1)$ and $-\tfrac{1}{\sqrt3}(1,-1,1)$. Evaluating,
\[f\Big(\tfrac{1}{\sqrt3}(1,-1,1)\Big) = \frac{3+1+2}{\sqrt3} = 2\sqrt3, \qquad f\Big({-\tfrac{1}{\sqrt3}}(1,-1,1)\Big) = -2\sqrt3.\]
Since $M$ is compact (closed and bounded in $\mathbb{R}^3$) and $f$ is continuous, these are the
only Lagrange candidates, so $f$ attains its \textbf{maximum} $2\sqrt3$ at $\tfrac{1}{\sqrt3}(1,-1,1)$
and its \textbf{minimum} $-2\sqrt3$ at $-\tfrac{1}{\sqrt3}(1,-1,1)$.
\end{exercisesolution}

\begin{exercisesolution}[Solution to \cref{ex:6.11}]
% Generator: Claude Sonnet 5 (Medium)
Minimizing the distance to $(0,0,1)$ is the same as minimizing
$f(x,y,z) := x^2+y^2+(z-1)^2$ subject to $g(x,y,z) := z-x^2+y^2 = 0$. Since $f$ is a distance
squared to a fixed point, its minimum over the (closed) constraint set is attained. Lagrange's
condition $\nabla f = \lambda \nabla g$ reads
\[2x = -2\lambda x, \qquad 2y = 2\lambda y, \qquad 2(z-1) = \lambda,\]
i.e.\ $x(1+\lambda) = 0$ and $y(1-\lambda) = 0$.

\textbf{Case $x=y=0$:} the constraint gives $z=0$, so $(0,0,0)$ with $d^2 = 1$.

\textbf{Case $x \neq 0$ (so $\lambda=-1$):} then $2(z-1)=-1 \implies z = \tfrac12$, and $y(1-\lambda)=2y=0
\implies y=0$. The constraint $z=x^2-y^2$ gives $x^2=\tfrac12$, i.e.\ $x=\pm\tfrac{1}{\sqrt2}$,
giving $\big(\pm\tfrac{1}{\sqrt2},0,\tfrac12\big)$ with $d^2 = \tfrac12+0+\tfrac14 = \tfrac34$.

\textbf{Case $y \neq 0$ (so $\lambda=1$):} then $2(z-1)=1 \implies z=\tfrac32$, and $x(1+\lambda)=2x=0
\implies x=0$; the constraint becomes $-y^2 = \tfrac32$, which has no real solution.

Comparing $d^2=1$ against $d^2=\tfrac34$, the \textbf{closest points} are
$\big(\pm\tfrac{1}{\sqrt2},0,\tfrac12\big)$, at minimal distance $\tfrac{\sqrt3}{2}$.
\end{exercisesolution}

